\documentclass[11pt,a4paper]{article}
\usepackage[T1]{fontenc}
\usepackage[utf8]{inputenc}
\usepackage{lmodern,amsmath,amssymb}
\usepackage[margin=25mm]{geometry}
\usepackage[hidelinks]{hyperref}
\hypersetup{pdftitle={A conservative receiver and the origin of a correlation scale},pdfauthor={navstokgap research working notes}}
\newcommand{\tightlist}{\setlength{\itemsep}{0pt}\setlength{\parskip}{0pt}}
\setlength{\emergencystretch}{3em}
\setcounter{secnumdepth}{0}
\begin{document}
\hypertarget{a-conservative-receiver-and-the-origin-of-a-correlation-scale}{%
\section{A conservative receiver and the origin of a correlation
scale}\label{a-conservative-receiver-and-the-origin-of-a-correlation-scale}}

A finite harmonic receiver has a vanishing long-window action observable
at fixed centre velocity. An increasing periodic chain instead has
plateau \(\Theta\sqrt{m/k}\) when its size tends to infinity first, with
mode energy \(\Theta\) and spring stiffness \(k\). A common hard
particle-speed ceiling on the stated phase preparation forces this
plateau to close. The comparison locates the positive scale in the
low-frequency weights and energy preparation.

\hypertarget{hamiltonian-and-invariant-preparation}{%
\subsection{1. Hamiltonian and invariant
preparation}\label{hamiltonian-and-invariant-preparation}}

Let \(x,p\in\mathbb R^n\), \(n\ge2\),
\(M=\operatorname{diag}(m_1,\ldots,m_n)\) with \(m_i>0\), and let \(K\)
be a connected spring-network Laplacian. Thus \(K=K^T\ge0\) and
\(\ker K=\operatorname{span}\{\mathbf1\}\). The Hamiltonian and
equations are

\[
\mathcal H=\tfrac12p^TM^{-1}p+\tfrac12x^TKx,\qquad
\dot x=M^{-1}p,\quad\dot p=-Kx.
\]

Energy and total momentum are conserved. These linear equations give
global smooth trajectories. Write \(D=M^{-1/2}KM^{-1/2}\) and choose an
orthonormal eigenbasis \(e_0,e_1,\ldots,e_{n-1}\), with
\(e_0=M^{1/2}\mathbf1/\sqrt{M_{\rm tot}}\) and positive eigenvalues
\(\omega_j^2\) for \(j\ge1\). Define canonical coordinates
\(q_j=e_j^TM^{1/2}x\), \(P_j=e_j^TM^{-1/2}p\).

Fix centre velocity \(V_0=0\) and centre position. Assign internal mode
energies \(E_j\ge0\), and independent phases \(\phi_j\) uniform on
\([0,2\pi)\). Then

\[
q_j(t)=\frac{\sqrt{2E_j}}{\omega_j}\cos(\omega_jt+\phi_j),\qquad
P_j(t)=-\sqrt{2E_j}\sin(\omega_jt+\phi_j).
\]

This is an invariant probability ensemble on the internal phase tori,
even when frequencies are commensurate. It is a preparation, not a
mixing premise. The total energy is \(E=\sum_jE_j\), so
\(|\dot x_i|\le\sqrt{2E/m_i}\) for every trajectory. Choosing
\(E<\tfrac12(\min_i m_i)c^2\) puts all particles below \(c\) in this
frame. These are Newtonian springs with a speed-bounded preparation;
propagation of signals and Lorentz covariance are separate model
requirements.

\hypertarget{exact-action-observable-and-its-limits-c047}{%
\subsection{2. Exact action observable and its limits
(C047)}\label{exact-action-observable-and-its-limits-c047}}

For tagged particle \(i\), let \(b_{ij}=e_{ji}/\sqrt{m_i}\), where
\(e_{ji}\) denotes component \(i\) of vector \(e_j\). Then
\(v_i(t)=\sum_{j\ge1}b_{ij}P_j(t)\) has zero mean and covariance

\[
C_i(t)=\sum_{j\ge1}w_{ij}\cos(\omega_jt),\qquad
w_{ij}=b_{ij}^2E_j\ge0.
\]

Phase averaging gives this identity: distinct modes have zero cross
covariance, and one phase has
\(\mathbb E[P_j(t)P_j(0)]=E_j\cos(\omega_jt)\). Integrating the
covariance over the displacement square yields, for \(\Delta>0\),

\[
\mathsf h_i(\Delta)=\frac{m_i}{\Delta}
\operatorname{Var}[x_i(t+\Delta)-x_i(t)]
=\frac{2m_i}{\Delta}\sum_{j\ge1}
\frac{w_{ij}}{\omega_j^2}[1-\cos(\omega_j\Delta)].
\]

The observable has action units and is independent of preparation time
\(t\). Set \(B_i=\sum_jw_{ij}/\omega_j^2\), with units of length
squared. Then

\[
0\le\mathsf h_i(\Delta)\le\frac{4m_iB_i}{\Delta}\longrightarrow0
\quad(\Delta\to\infty).
\]

At short windows, \(\mathsf h_i(\Delta)=m_iC_i(0)\Delta+O(\Delta^3)\).
The finite cosine correlation has persistent memory; its ordinary
improper Green--Kubo integral need not exist. The finite-window formula
above supplies the relevant limit directly, without replacing the
Hamiltonian dynamics by a reversible dissipative generator.

\hypertarget{centre-motion-and-the-two-body-receiver-c047}{%
\subsection{3. Centre motion and the two-body receiver
(C047)}\label{centre-motion-and-the-two-body-receiver-c047}}

An independent random constant centre velocity with variance
\(\sigma_0^2\) adds \(m_i\sigma_0^2\Delta\) to \(\mathsf h_i\). This
follows from \(x_i(t+\Delta)-x_i(t)=V_0\Delta+\) internal displacement.
The velocity ensemble is stationary; for nonzero centre motion the
absolute-position ensemble need not be. A deterministic centre velocity
contributes zero to the variance and can be removed by a change of
frame.

For two particles of masses \(m,M_r\) coupled by a spring \(k>0\), put
\(\mu=mM_r/(m+M_r)\) and \(\omega^2=k/\mu\). At zero total momentum,
\(x_1=R+[M_r/(m+M_r)]r\), and assign relative energy \(E\) with uniform
phase. The tagged response is

\[
\mathsf h_1(\Delta)=
\frac{2mE}{k\Delta}\left(\frac{M_r}{m+M_r}\right)^2
[1-\cos(\omega\Delta)].
\]

The spring is a genuine momentum receiver: \(m\ddot x_1=-kr\) and
\(M_r\ddot x_2=kr\). Its reversible exchange produces the oscillatory
response rather than the exponential memory law supplied by A02's
refreshed bath. Scaling every \(E_j\) by \(a^2\), \(0<a\le1\), scales
the entire response by \(a^2\) while preserving the equations,
conservation laws and speed margin.

\hypertarget{what-the-large-receiver-test-must-change-c047}{%
\subsection{4. What the large-receiver test must change
(C047)}\label{what-the-large-receiver-test-must-change-c047}}

For a sequence of receivers and windows \(\Delta_N\to\infty\), the bound
\(\mathsf h_{i,N}(\Delta_N)\le4m_{i,N}B_{i,N}/\Delta_N\) shows that a
positive limiting response requires this upper bound to stay positive.
In particular, uniform control of \(m_{i,N}B_{i,N}\) forces a zero
limit. Soft modes can enlarge \(B_{i,N}\), but their weights depend on
the measured coordinate and the preparation. This is the next selection
test.

The finite model yields no physical-time attraction from the invariant
preparation: the observable already has no \(t\) dependence. Next
specify an increasing spring network, its energy allocation and tagged
spectral measure, then compare taking receiver size and observation
duration to infinity in opposite orders. Derive any surviving constant
from these mechanical inputs. Mesh refinement of the same smooth
trajectories is a third, distinct limit.

\hypertarget{an-increasing-periodic-chain-c048}{%
\subsection{5. An increasing periodic chain
(C048)}\label{an-increasing-periodic-chain-c048}}

Take odd \(N\ge3\) equal masses \(m\) on a ring, with spring energy
\(\frac{k}{2}\sum_{i=0}^{N-1}(x_{i+1}-x_i)^2\) (indices modulo \(N\)).
Fix the centre position and momentum. Give every real internal normal
mode energy \(\Theta>0\) and independent uniform phase. Here \(\Theta\)
is a preparation energy, and total energy is \((N-1)\Theta\). The mode
frequencies and tagged covariance are

\[
\omega_j=\Omega\sin(\pi j/N),\quad \Omega=2\sqrt{k/m},\qquad
C_N(t)=\frac{\Theta}{mN}\sum_{j=1}^{N-1}\cos(\omega_jt).
\]

The real sine/cosine eigenvectors have equal mode energies, so their
paired squared components sum to \(2/N\) at each site. Applying §2 gives

\[
h_N(\Delta)=\frac{2\Theta}{N\Delta}\sum_{j=1}^{N-1}
\frac{1-\cos(\omega_j\Delta)}{\omega_j^2}.
\]

For each fixed \(\Delta\), this Riemann sum converges to

\[
h_\infty(\Delta)=\frac{2m}{\Delta}\int_0^\Omega
\frac{1-\cos(\omega\Delta)}{\omega^2}\rho(\omega)\,d\omega,
\quad \rho(\omega)=\frac{2\Theta}{\pi m\sqrt{\Omega^2-\omega^2}}.
\]

This is a statement about limits of tagged covariances and increment
variances; it does not require a stationary absolute-position law for
the infinite chain. The density is integrable at its upper edge and
continuous at zero. Splitting at any fixed \(0<a<\Omega\), the
contribution from \([a,\Omega]\) is \(O(1/\Delta)\). In the lower
interval substitute \(y=\omega\Delta\). Dominated convergence, using
\((1-\cos y)/y^2\in L^1(0,\infty)\) and bounded \(\rho\) on \([0,a]\),
gives

\[
\lim_{\Delta\to\infty}h_\infty(\Delta)
=\pi m\rho(0)=\frac{2\Theta}{\Omega}=\Theta\sqrt{m/k}.
\]

The integral \(\int_0^\infty(1-\cos y)y^{-2}dy=\pi/2\) follows by
integration by parts and the Dirichlet sine integral. Thus

\[
\lim_{\Delta\to\infty}\lim_{N\to\infty}h_N(\Delta)=2\Theta/\Omega,
\qquad
\lim_{N\to\infty}\lim_{\Delta\to\infty}h_N(\Delta)=0.
\]

A joint regime is also explicit. Write
\(f(z)=2(1-\cos z)/z^2=2\int_0^1(1-r)\cos(zr)\,dr\) with \(f(0)=1\).
Then \(|f'|\le1/3\). The integrand
\(\Theta\Delta f(\Omega\Delta\sin\pi x)\) is Lipschitz with constant at
most \(\pi\Theta\Omega\Delta^2/3\). A left Riemann sum error is at most
that constant divided by \(2N\); removing its zero mode contributes
\(\Theta\Delta/N\). Therefore

\[
|h_N(\Delta)-h_\infty(\Delta)|
\le\frac{\pi\Theta\Omega\Delta^2}{6N}+\frac{\Theta\Delta}{N}.
\]

At fixed \(m,k,\Theta\), \(\Omega\Delta_N\to\infty\) and
\((\Omega\Delta_N)^2/N\to0\) give the positive joint limit
\(2\Theta/\Omega\). The action scale is selected by the input energy and
spring timescale.

\hypertarget{energy-and-a-common-speed-ceiling-c049}{%
\subsection{6. Energy and a common speed ceiling
(C049)}\label{energy-and-a-common-speed-ceiling-c049}}

The positive-plateau preparation has extensive total energy. Its phase
support also contains increasingly large tagged velocities. At site zero
use the standard real Fourier basis. Only the \((N-1)/2\) cosine modes
contribute; each has maximal velocity contribution
\(2\sqrt{\Theta/(mN)}\). Phases can align, so the exact support maximum
is

\[
\sup_{\phi}|v_0|=(N-1)\sqrt{\Theta/(mN)}.
\]

Every neighbourhood of the maximizing phases has positive preparation
measure. A common ceiling \(|v_i|\le c\) for all supported trajectories
thus requires \(\Theta_N\le mc^2N/(N-1)^2=O(1/N)\). This is a necessary
condition from one site; no sufficiency for all sites is inferred from
that maximum.

Both bounded total energy and this necessary speed condition close the
response uniformly in the window. Indeed, \(1-\cos z\le z^2/2\) gives
\(h_N(\Delta)\le\Theta_N\Delta\). The identity
\(\sum_{j=1}^{N-1}\csc^2(\pi j/N)=(N^2-1)/3\) gives

\[
h_N(\Delta)\le
\frac{4\Theta_N(N^2-1)}{3N\Omega^2\Delta},\qquad
\sup_{\Delta>0}h_N(\Delta)
\le\frac{2\Theta_N}{\Omega}\sqrt{\frac{N^2-1}{3N}}.
\]

The trigonometric identity follows by differentiating
\(\sum_{j=0}^{N-1}\cot(x+\pi j/N)=N\cot(Nx)\) and subtracting the
\(j=0\) singularity as \(x\to0\). The last inequality uses
\(\min(A\Delta,B/\Delta)\le\sqrt{AB}\). For \(\Theta_N=O(1/N)\) the
uniform bound tends to zero. Thus this explicit positive reservoir
plateau uses a different preparation class from a family with a common
strict particle-speed ceiling. Its dependence on \(\Theta\) also leaves
preparation-independent action selection open.

\hypertarget{source-input-and-reproduction}{%
\subsection{7. Source input and
reproduction}\label{source-input-and-reproduction}}

Ford--Kac--Mazur, printed p.~505, (2)--(5), supplies the matrix harmonic
propagator. Their canonical preparation excludes zero eigenvalues at
that stage; ours fixes the centre and uses compact fixed-energy phase
tori. Zwanzig, p.~219, (21)--(24), supplies a complementary
reduced-description route: eliminate harmonic bath coordinates to obtain
a cosine memory kernel. That friction kernel differs from the tagged
velocity covariance calculated here. Its continuum replacement must be
propagated through the tracer equation before identifying a displacement
coefficient.

\href{../references/batches/B23.md}{B23} records the established
ingredients, derived finite-network consequence and bounded reading
coverage. Historical modal, two-body and three-body checks accompany
C047's proof. C047 has coordinator proof review and a sequential
Luna-low prior-art audit. Its limits concern the stated observable and
preparation class.

\href{../references/batches/B24.md}{B24} audits §§5--6 against
Ford--Kac--Mazur's cyclic Fourier construction and its
finite-sum/infinite-chain distinction, printed p.~506, (12)--(24). The
tagged density, plateau and speed-support bounds are derived here with
fixed-energy phases. Six exact identities, four finite ring
spectrum/cosecant comparisons and four Riemann-bound tests were recorded
before the hard no-Python-verification rule; the handoff records that
history. The new uncommitted verification script and output were
withdrawn under the rule. Current mathematical verification uses the
written proofs and source review. The next mechanism must retain bounded
velocities while supplying low-frequency response; its preparation
dependence remains a separate selection test.

\hypertarget{what-information-survives-a-cut}{%
\subsection{8. What information survives a
cut?}\label{what-information-survives-a-cut}}

The fixed-energy two-body receiver gives an exact cut-state countertest:
position-only stationary conditional kernels fail composition, while
retaining momentum restores it. Independently refreshing direction at
every inserted cut preserves one-time position marginals but freezes
terminal position as the mesh shrinks. These are C062--C063; the full
maps and proof are in \href{classical-cut-state.md}{the cut-state note},
with the \href{../references/batches/B33.md}{B33 audit}.

With \(q=x-X\), \(\mu=mM/(m+M)\), \(\omega=\sqrt{k/\mu}\) and
\(A=\sqrt{2E/k}\), uniform phase on the energy ellipse gives the
conditional one-segment positions

\[q'=q\cos\omega t\ \pm\sqrt{A^2-q^2}\sin\omega t\]

with equal weights. Thus the full mean after \(s+t\) is
\(q\cos\omega(s+t)\), whereas independent conditional composition gives
\(q\cos\omega s\cos\omega t\). The missing interface variable is
\(p_s=\mu\dot q_s\) in
\(q_{s+t}=q_s\cos\omega t+p_s\sin\omega t/(\mu\omega)\). It has momentum
units. Keeping this full internal state composes exactly; marginalizing
the true joint path law also preserves all cut restrictions.

For \(N\) independent resets at intervals \(T/N\), direct iteration
gives

\[\mathbb E Q_N=q[\cos(\omega T/N)]^N,\] \[\mathbb E Q_N^2=\frac{A^2}{2}
 +\left(q^2-\frac{A^2}{2}\right)[\cos(2\omega T/N)]^N.\]

The first limit is \(q\) and the second is \(q^2\), proving terminal
conditional \(L^2\) convergence to the starting position. Under the
original stationary preparation the reset two-time correlation tends to
\(A^2/2\), while the true value is \((A^2/2)\cos\omega T\). A cut that
only observes the motion performs no reset and leaves the true
correlation unchanged.

The normalized orbital action is \(J=(2\pi)^{-1}\oint p\,dq=E/\omega\):
integrating \(\mu\omega A^2\sin^2\phi\) over a cycle proves the formula.
Uniform-phase covariance also has square-root determinant \(E/\omega\).
Cooling \(E\downarrow0\) at fixed Hamiltonian closes both, with exact
phase-state composition intact. R04 extends this test to a third body,
as follows.

\hypertarget{a-third-body-exact-memory-and-recoverable-history}{%
\subsection{9. A third body: exact memory and recoverable
history}\label{a-third-body-exact-memory-and-recoverable-history}}

Retaining tagged momentum leaves receiver memory in the equal-mass open
three-body chain. At zero centre position and momentum, let \(x=x_1\),
\(y=x_2-x_3\), \(\mu=3m/2\), \(\nu=m/2\), \(P=\mu\dot x\) and
\(Q=\nu\dot y\). Physical tagged momentum is \(p_1=2P/3\). The internal
Hamiltonian is

\[H=\frac{P^2}{2\mu}+\frac{Q^2}{2\nu}
+\frac{9k}{8}x^2-\frac{3k}{4}xy+\frac{5k}{8}y^2.\]

For fixed microcanonical energy \(E>0\), eliminate \(y,Q\) exactly. With
\(g=3k/4\), \(d=5k/4\) and \(\Omega^2=5k/(2m)\),

\[\mu\ddot x+\frac{9k}{5}x
+\int_0^t\frac{9k}{20}\cos\Omega(t-s)\dot x(s)\,ds=F_0(t),\]

\[F_0(t)=g\left(y_0-\frac gd x_0\right)\cos\Omega t
+\frac{gQ_0}{\nu\Omega}\sin\Omega t.\]

These inherited quadratures carry the receiver information across a cut.
The kernel has stiffness units; its convolution is a force. The full
proof in \href{three-body-cut-memory.md}{R04} also recovers the hidden
pair from tagged phase at two sufficiently close times. The relevant
block determinant is \(g^2\delta^4/(12\mu\nu)+O(\delta^6)\); inverse
entries grow up to \(\delta^{-3}\). Exact recovery and finite-precision
physical readout therefore pose different requirements.

The internal frequencies are \(\sqrt{k/m}\) and \(\sqrt{3k/m}\), and
modal actions are \(E_j/\omega_j\). Cooling scales the actions to zero
while keeping the memory kernel fixed. C064--C065 and
\href{../references/batches/B34.md}{B34} record this derived
specialization of harmonic-bath elimination and observability. R05 next
specifies the measurement interaction and tests its precision,
disturbance and resource scaling.
\end{document}
